% Conjectural application.  The proved Hasse-jet results do not depend on it.

\section{Tate-dual obstruction classes}

Put $P=p(p-1)$ and let $i=np+i'$ be in the nondegenerate band. Define
\[
  i^\dagger=p^2+1-k-i
   =(p-2-n)p+(2p+1-k-i').
\]
Then
\[
 i+i^\dagger=p^2+1-k,
 \qquad
 (k+2i)+(k+2i^\dagger)=2p^2+2.
\]
Moreover, if $c(x)=x^2+(k-1)x$, then
$c(2p+1-k-i')\equiv c(i')\pmod p$, while
\[
 (p-2-n)(p-1-n)\equiv(n+1)(n+2)\pmod p.
\]

The lower Bockstein class $\beta_w(g)$ and its lift-independence are given
by Lemma~\ref{lem:bockstein}. In particular,
\[
  \beta_w(g)=0
  \quad\Longleftrightarrow\quad
  \omega_{p^2}(g)\le w
\]
for an admissible candidate weight $w$.

The perfect Hasse-quotient pairing and theta adjointness are proved in
Propositions~\ref{prop:hasse-perfect-pairing} and
\ref{prop:theta-adjointness}. Put
\[
 w_-=k+2i-P,\qquad w_+=k+2i^\dagger+P.
\]
Proposition~\ref{prop:geometric-bockstein} places the two actual
weight-lowering obstructions in
\[
 Q_{r_-}\quad\hbox{and}\quad Q_{r_+},
 \qquad
 r_-=w_-+P=k+2i,
 \quad
 r_+=w_++P=k+2i^\dagger+2P.
\]
Their weights satisfy
\[
 r_-+r_+=P+2+p(3p-1).
\]
Theorem~\ref{thm:natural-bockstein-pairing} uses a horizontal unit of
weight $p(3p-1)$ to give a perfect pairing
\[
 [\ ,\ ]^{\mathrm{nat}}:
 Q_{r_-}\times Q_{r_+}\longrightarrow\mathbb F_p
\]
under which theta is skew-adjoint.

\begin{conjecture}[Cross-wedge Bockstein compatibility]
Assume that $f$ is ordinary at $p$, $4\le k<p-1$, and that both $i$ and
$i^\dagger$ lie in the nondegenerate band. If
\[
 D(i)=0,\qquad D(i^\dagger)=2,
\]
then the two geometric Bockstein classes
\[
 \Gamma_{w_-}(\theta^if)\in Q_{r_-},
 \qquad
 \Gamma_{w_+}(\theta^{i^\dagger}f)\in Q_{r_+}
\]
represent one another under the natural perfect pairing, up to the explicit
nonzero scalar in the $D=2$ block formula. Consequently
\[
 \beta_{k+2i-P}(\theta^i f)=0
 \quad\Longleftrightarrow\quad
 \beta_{k+2i^\dagger+P}(\theta^{i^\dagger}f)=0.
\]
\end{conjecture}

\begin{conjecture}[Exact $D=0$ shifted-conic law]
Under the hypotheses of the preceding conjecture,
\[
 \delta(\theta^i f)=-1
 \quad\Longleftrightarrow\quad
 i'^2+(k-1)i'-(n+1)(n+2)\equiv0\pmod p.
\]
Otherwise $\delta(\theta^i f)=0$.
\end{conjecture}

\begin{proof}[Deduction of the second conjecture from the first]
For the dual index, the proved $D=2$ polynomial is
\[
 (i'^\dagger)^2+(k-1)i'^\dagger
   -n^\dagger(n^\dagger+1).
\]
The two elementary congruences above transform this into
\[
 i'^2+(k-1)i'-(n+1)(n+2).
\]
The $D=2$ theorem says that the complementary obstruction vanishes exactly
on this conic. Cross-wedge Tate duality transfers the vanishing to $i$.
\end{proof}

\begin{conjecture}[Corrected global duality]
Let $e(i)=D(i)-\delta(\theta^i f)$. Away from the central $D=1$ plateau,
\[
 e(i)=e(i^\dagger).
\]
On that plateau the difference is a rank-one boundary functional supported
on
\[
 i'^2+(k-1)i'-(n+1)^2\equiv0\pmod p.
\]
\end{conjecture}
